\clearpage\thispagestyle{empty}
\pdfbookmark[0]{Quellen zu Gauss an Bode (22. August 1812)}{bode1812aug22-sources}
\begin{center}\textbf{Quellen zu Gauss an Bode (22. August 1812)}\end{center}
\noindent Diese KI-Übertragung ins heutige Deutsch folgt dem Brief von Gauß vom 22. August 1812 in den \href{https://gdz.sub.uni-goettingen.de/id/PPN236007467}{\textit{Werken}, Band~VI} (Göttingen, 1874), S.~357--360. Der Erstdruck erschien im \href{https://api.digitale-sammlungen.de/iiif/image/v2/bsb10538344_00005/full/1600,/0/default.jpg}{\textit{Astronomisches Jahrbuch für das Jahr 1815}}, herausgegeben von Johann Elert Bode (Berlin, 1812), S.~\href{https://api.digitale-sammlungen.de/iiif/image/v2/bsb10538344_00253/full/1600,/0/default.jpg}{245}, \href{https://api.digitale-sammlungen.de/iiif/image/v2/bsb10538344_00254/full/1600,/0/default.jpg}{246}, \href{https://api.digitale-sammlungen.de/iiif/image/v2/bsb10538344_00255/full/1600,/0/default.jpg}{247}, \href{https://api.digitale-sammlungen.de/iiif/image/v2/bsb10538344_00256/full/1600,/0/default.jpg}{248}, \href{https://api.digitale-sammlungen.de/iiif/image/v2/bsb10538344_00257/full/1600,/0/default.jpg}{249}. Das digitalisierte Exemplar stammt aus der Bayerischen Staatsbibliothek.
\par\smallskip\noindent Im Brief der Werkausgabe steht bei den Juno-Beobachtungen Januar und beim Vergleich Juni. Diese Monatsangaben finden sich auch im Erstdruck. Dort stehen auf S.~248 bei der Epoche das Jahr 1811 und bei der letzten Deklinationsdifferenz ein Dezimalkomma. In der Werkausgabe fehlt bei der Epoche die Jahresangabe, und die Differenz lautet $-45\;0$. Das Jahrbuch enthält auf S.~249 drei mit B. gezeichnete Anmerkungen, darunter einen Verweis auf die im Jahrbuch folgende Juno-Ephemeride. Diese Anmerkungen gehören zum Erstdruck. KI-Übertragung und KI-Satz.
\clearpage
